% Options for packages loaded elsewhere
% Options for packages loaded elsewhere
\PassOptionsToPackage{unicode}{hyperref}
\PassOptionsToPackage{hyphens}{url}
\PassOptionsToPackage{dvipsnames,svgnames,x11names}{xcolor}
% Code above from https://github.com/quarto-dev/quarto-cli/blob/main/src/resources/formats/pdf/pandoc/passoptions.latex
% Customization below
\PassOptionsToPackage{babel}{csquotes}
%
\documentclass[
  french,
  11pt,
]{stylisharticle}
\usepackage{xcolor}
\usepackage{amsmath,amssymb}
\setcounter{secnumdepth}{2}
\usepackage{iftex}
\ifPDFTeX
  \usepackage[T1]{fontenc}
  \usepackage[utf8]{inputenc}
  \usepackage{textcomp} % provide euro and other symbols
\else % if luatex or xetex
  \usepackage{unicode-math} % this also loads fontspec
  \defaultfontfeatures{Scale=MatchLowercase}
  \defaultfontfeatures[\rmfamily]{Ligatures=TeX,Scale=1}
\fi
\usepackage{lmodern}
\ifPDFTeX\else
  % xetex/luatex font selection
\fi
% Use upquote if available, for straight quotes in verbatim environments
\IfFileExists{upquote.sty}{\usepackage{upquote}}{}
\IfFileExists{microtype.sty}{% use microtype if available
  \usepackage[]{microtype}
  \UseMicrotypeSet[protrusion]{basicmath} % disable protrusion for tt fonts
}{}
\makeatletter
\@ifundefined{KOMAClassName}{% if non-KOMA class
  \IfFileExists{parskip.sty}{%
    \usepackage{parskip}
  }{% else
    \setlength{\parindent}{0pt}
    \setlength{\parskip}{6pt plus 2pt minus 1pt}}
}{% if KOMA class
  \KOMAoptions{parskip=half}}
\makeatother
% Make \paragraph and \subparagraph free-standing
\makeatletter
\ifx\paragraph\undefined\else
  \let\oldparagraph\paragraph
  \renewcommand{\paragraph}{
    \@ifstar
      \xxxParagraphStar
      \xxxParagraphNoStar
  }
  \newcommand{\xxxParagraphStar}[1]{\oldparagraph*{#1}\mbox{}}
  \newcommand{\xxxParagraphNoStar}[1]{\oldparagraph{#1}\mbox{}}
\fi
\ifx\subparagraph\undefined\else
  \let\oldsubparagraph\subparagraph
  \renewcommand{\subparagraph}{
    \@ifstar
      \xxxSubParagraphStar
      \xxxSubParagraphNoStar
  }
  \newcommand{\xxxSubParagraphStar}[1]{\oldsubparagraph*{#1}\mbox{}}
  \newcommand{\xxxSubParagraphNoStar}[1]{\oldsubparagraph{#1}\mbox{}}
\fi
\makeatother


\usepackage{longtable,booktabs,array}
\newcounter{none} % for unnumbered tables
\usepackage{calc} % for calculating minipage widths
% Correct order of tables after \paragraph or \subparagraph
\usepackage{etoolbox}
\makeatletter
\patchcmd\longtable{\par}{\if@noskipsec\mbox{}\fi\par}{}{}
\makeatother
% Allow footnotes in longtable head/foot
\IfFileExists{footnotehyper.sty}{\usepackage{footnotehyper}}{\usepackage{footnote}}
\makesavenoteenv{longtable}
\usepackage{graphicx}
\makeatletter
\newsavebox\pandoc@box
\newcommand*\pandocbounded[1]{% scales image to fit in text height/width
  \sbox\pandoc@box{#1}%
  \Gscale@div\@tempa{\textheight}{\dimexpr\ht\pandoc@box+\dp\pandoc@box\relax}%
  \Gscale@div\@tempb{\linewidth}{\wd\pandoc@box}%
  \ifdim\@tempb\p@<\@tempa\p@\let\@tempa\@tempb\fi% select the smaller of both
  \ifdim\@tempa\p@<\p@\scalebox{\@tempa}{\usebox\pandoc@box}%
  \else\usebox{\pandoc@box}%
  \fi%
}
% Set default figure placement to htbp
\def\fps@figure{htbp}
\makeatother


% definitions for citeproc citations
\NewDocumentCommand\citeproctext{}{}
\NewDocumentCommand\citeproc{mm}{%
  \begingroup\def\citeproctext{#2}\cite{#1}\endgroup}
\makeatletter
 % allow citations to break across lines
 \let\@cite@ofmt\@firstofone
 % avoid brackets around text for \cite:
 \def\@biblabel#1{}
 \def\@cite#1#2{{#1\if@tempswa , #2\fi}}
\makeatother
\newlength{\cslhangindent}
\setlength{\cslhangindent}{1.5em}
\newlength{\csllabelwidth}
\setlength{\csllabelwidth}{3em}
\newenvironment{CSLReferences}[2] % #1 hanging-indent, #2 entry-spacing
 {\begin{list}{}{%
  \setlength{\itemindent}{0pt}
  \setlength{\leftmargin}{0pt}
  \setlength{\parsep}{0pt}
  % turn on hanging indent if param 1 is 1
  \ifodd #1
   \setlength{\leftmargin}{\cslhangindent}
   \setlength{\itemindent}{-1\cslhangindent}
  \fi
  % set entry spacing
  \setlength{\itemsep}{#2\baselineskip}}}
 {\end{list}}
\usepackage{calc}
\newcommand{\CSLBlock}[1]{\hfill\break\parbox[t]{\linewidth}{\strut\ignorespaces#1\strut}}
\newcommand{\CSLLeftMargin}[1]{\parbox[t]{\csllabelwidth}{\strut#1\strut}}
\newcommand{\CSLRightInline}[1]{\parbox[t]{\linewidth - \csllabelwidth}{\strut#1\strut}}
\newcommand{\CSLIndent}[1]{\hspace{\cslhangindent}#1}

\ifLuaTeX
\usepackage[bidi=basic,shorthands=off]{babel}
\else
\usepackage[bidi=default,shorthands=off]{babel}
\fi
\ifLuaTeX
  \usepackage{selnolig} % disable illegal ligatures
\fi


\setlength{\emergencystretch}{3em} % prevent overfull lines

\providecommand{\tightlist}{%
  \setlength{\itemsep}{0pt}\setlength{\parskip}{0pt}}



 

\usepackage[]{csquotes}

% -------------- Additional packages
\usepackage{orcidlink}  % Orcid logo and link to author page

% -------------- Font size
% R output: footnotesize
\makeatletter
\patchcmd{\@verbatim}
  {\verbatim@font}
  {\verbatim@font\footnotesize}
  {}{}
% References font size whether bibtex is used or not
\@ifundefined{bibfont}
  {%
    \AtBeginEnvironment{CSLReferences}{\footnotesize}
  }
  {%
    \renewcommand*{\bibfont}{\footnotesize}
  }%
\makeatother

% -------------- longtable 2 columns
% https://tex.stackexchange.com/questions/161431/how-to-solve-longtable-is-not-in-1-column-mode-error
\makeatletter
\let\oldlt\longtable
\let\endoldlt\endlongtable
\def\longtable{\@ifnextchar[\longtable@i \longtable@ii}
\def\longtable@i[#1]{
\begin{figure}[t]
\onecolumn
\begin{minipage}{0.5\textwidth}\scriptsize
\oldlt[#1]
}
\def\longtable@ii{
\begin{figure}[t]
\onecolumn
\begin{minipage}{0.5\textwidth}\scriptsize
\oldlt
}
\definecolor{red}{RGB}{255, 0, 0}
\definecolor{green}{RGB}{0, 255, 0}
\definecolor{blue}{RGB}{0, 0, 255}
\definecolor{grey}{RGB}{191, 191, 191}
\makeatletter
\@ifpackageloaded{caption}{}{\usepackage{caption}}
\AtBeginDocument{%
\ifdefined\contentsname
  \renewcommand*\contentsname{Table des matières}
\else
  \newcommand\contentsname{Table des matières}
\fi
\ifdefined\listfigurename
  \renewcommand*\listfigurename{Liste des Figures}
\else
  \newcommand\listfigurename{Liste des Figures}
\fi
\ifdefined\listtablename
  \renewcommand*\listtablename{Liste des Tables}
\else
  \newcommand\listtablename{Liste des Tables}
\fi
\ifdefined\figurename
  \renewcommand*\figurename{Figure}
\else
  \newcommand\figurename{Figure}
\fi
\ifdefined\tablename
  \renewcommand*\tablename{Table}
\else
  \newcommand\tablename{Table}
\fi
}
\@ifpackageloaded{float}{}{\usepackage{float}}
\floatstyle{ruled}
\@ifundefined{c@chapter}{\newfloat{codelisting}{h}{lop}}{\newfloat{codelisting}{h}{lop}[chapter]}
\floatname{codelisting}{Listing}
\newcommand*\listoflistings{\listof{codelisting}{Liste des Listings}}
\makeatother
\makeatletter
\makeatother
\makeatletter
\@ifpackageloaded{caption}{}{\usepackage{caption}}
\@ifpackageloaded{subcaption}{}{\usepackage{subcaption}}
\makeatother
\usepackage{bookmark}
\IfFileExists{xurl.sty}{\usepackage{xurl}}{} % add URL line breaks if available
\urlstyle{same}
\hypersetup{
  pdftitle={Demo Stylish Article},
  pdfauthor={Félix Romiguière},
  pdflang={fr-FR},
  pdfkeywords={template, demo},
  colorlinks=true,
  linkcolor={blue},
  filecolor={Maroon},
  citecolor={Blue},
  urlcolor={blue},
  pdfcreator={LaTeX via pandoc}}


% Set up hyperref, after loading it
\hypersetup{
  linkcolor=black,
  citecolor=black,
  colorlinks=true
}

% Journal information
  \JournalInfo{Publication reference}

% Additional notes (e.g. copyright, DOI, review/research article)
  \Archive{DOI: xxx/xx}

% Article title
\PaperTitle{Demo Stylish Article}

% Authors
\Authors{
      {Félix Romiguière}%
    \textsuperscript{%
      ~\orcidlink{0000-0000-0000-0000}%
      1%
      %
    }%
    }

% Affiliations
    \affiliation{%
      \textsuperscript{1}%
      Université de Montpellier%
      , Scientific Department%
      %
      , Somewhere%
      , Montpellier%
      %
      , France%
      , 9999%
      %
    }

% Corresponding author
  %

% Keywords
\Keywords{template, demo}
\newcommand{\keywordname}{Keywords} % Defines the keywords heading name

%	Abstract
\Abstract{
  Document pour le début de rédaction de l'article sur Tahiti Soundscape
}
\begin{document}
% Makes all text pages the same height
\flushbottom

% Print the title and abstract box
\maketitle

% Print the contents section
\tableofcontents

% Remove page numbering from the first page
\thispagestyle{empty}

\section{Introduction}\label{introduction}

\emph{Oceanic island urbanisation and non -native species invasion}

Oceanic island are small sized with unique biodiversity and simplified
communities shaped by human influence. They so stands out as unique
ecosystems useful to study the impacts of human activities (Whittaker et
al. 2017; Russell et Kueffer 2019). Oceanic island tends to be more
sensitive to species introductions and extensive anthropogenic land-used
modification (Whittaker et al. 2017; Spatz et al. 2017; Russell et
Kueffer 2019). For island in particular, avian species have been among
the most concerned taxa, having their community assemblages shaped by
both invasion and extensive urbanisation processes (Boyer et Jetz 2014;
Spatz et al. 2017; Rogers et al. 2021).

\emph{Interspecific competition and homogenization in islands}

Non-native birds are issued from specific families (e.g.,
Galliniformes,Columbidae, Fringillidae, Passeridae, Estrildidae and
Sturnidae ..) mainly introduced for hunting and aesthetic purposes
(Blackburn et Duncan 2001; Dyer et al. 2017). These species show
peculiar traits adapted to urban and agricultural landscapes, whereas
native birds from oceanic island tends to mostly include species
occuring in forests (Soares et al. 2021). Non-native species may
outcompete native for essential ressources, and can affect population
dynamic for native species (Davis 2003). As shown with the introduction
of the Japanese white eye in Hawai'i (Blackburn et al. 2009).
Interspecific competition can stand as a strong driver for native
non-native species assemblages (McGill et al. 2006). Combined effect of
the growing loss of native ecosystems and continued introduction of
birds species, especially of forest-dwelling birds, intensify non-native
competition threats (Sax et al. 2007). Moreover, the introduction of
functionnaly similar species, promotes negative impact through
functionnal homogenisation of island bird assemblages (Sobral et al.
2016; Sayol et al. 2021; Soares et al. 2022).

\emph{General context of Soundscape and ecoacoustic, acoustic indices}

Advance in soundscape ecology, enabled better understanding of structure
and dynamic of the sonic environments (Pijanowski et al. 2011). The
soundscape incorpore all components, classified according to the sound
sources: geophony (climate and geography), biophony (all living
organisms), and anthrophony (human activities) (Pijanowski et al. 2011).
Complexity of these acoustic emissions can be characterise via acoustic
indices that creates a numeric summary of the soundscape (Farina et
Pieretti 2012). Ecological meaning is inferred by relating acoustic
indices to diversity metrics (Towsey et al. 2014; Eldridge et al. 2018;
Dröge et al. 2021), habitat conditions (Fuller et al. 2015;
Sánchez-Giraldo et al. 2021; Rudge et al. 2026) or temporal dynamics
(Gage et Axel 2014). They can also provide indicators on soundscape
disturbances (Villanueva-Rivera et al. 2011).

\emph{Implication for oceanic islands}

In oceanic islands, ecological complexity is threathened by human
pressure (urbanisation, invasion). Soundscape studies offers a good tool
to monitor invasions, extinction, and assemblage turnover (Farina et
Pieretti 2012; Ribeiro et al. 2022). Such example exist with the study
of an invasive ant leading to impoverishment of the soundscape in New
Caledonian forests (Gasc et al. 2018). Species assemblages interact
acoustically and reflect a dynamic processes called an acoustic
community (Lellouch et al. 2014; Farina et James 2016). In oceanic
islands, acoustic communities are modified by non-native species
introduction, having strong implications for native species acoutic
habitat and their soundtope quality modification (Farina 2014; Merchant
et al. 2015). The soundtope represent the acoustic pattern created by
the distribution of biophonies (Farina 2014). Whereas the acoustic
habitat incorpore the soundtope and add the environmental acoustic
conditions in which the species assemblage is leaving (Merchant et al.
2015; Farina et James 2016).

\emph{Acoustic implication}

Recent studies have pointed out the potential importance of invasive
species as underappreciated source of novel noise, that could alter
acoustic communication among native species (Hopkins et al. 2022) (Hunt
et al. 2024). In the context of oceanic island, diminished bird richness
reflect fewer vocalizing species, resulting in less saturated
soundscapes than mainland (Robert et al. 2019). So, acoustic
interferences are anticipated to lessened on islands. Yet, introduction
of non-native species are numerous and is expected to continue to
increase in islands (Seebens et al. 2017). Compared to native species,
loudness and high occurrence rate of vocalization can be characteristic
for invasive species (Hopkins et al. 2022). For example, over the course
of one year, invasive Pekin Robins' (\emph{Leiothrix lutea}) calls made
up 37\% of all bird songs in a shrubland in their invaded range in
Europe (Farina et al. 2013). Likewise, calls of invasive birds such as
Common Mynas (\emph{Acridotheres tristis}) seem to dominate the urban
environments they have invaded (personal information from (Hopkins et
al. 2022)). Plus, timing of calls between native and invasive are very
similar, and may creates overlap, important for acoustic competition
(Farina et al. 2013; Bleach et al. 2015)

This novel biotic noises are quite attuned to mask vital information
encoded in acoustic signals, and may impact frequencies, reduce
perception and discrimination of the latter (Lohr et al. 2003). This
masking effect can negatively affect communication, movement, vigilance,
mating and foraging (Shannon et al. 2016). Affecting individual survival
and reproductive success (Halfwerk et al. 2011; Medeiros et al. 2017).
Therefore, the calls of non-native species, can bring a significant
avenue for acoustic competition with native species and may dominates
insular soundscape which favored it's homogenization (Farina et al.
2013). Still, evidence in oceanic island are scarce and need further
investigation (Hopkins et al. 2022; Hunt et al. 2024).

\emph{Problematisation}

Here, we sought to investigate the variation in saturation, intensity,
complexity of signal emissions at the acoustic habitat and soundtope
level of the soundscape, in relation to native/non-native bird
assemblage and habitat factors, across the habitat gradients Tahiti. We
predicted that urbanisation gradients drives relevant acoustic indices
results for the acoustic habitat level (Fairbrass et al. 2017). We also
predicted that differences in bird assemblage would drive the variation
at the soundtope level, and that non-native species in these assemblages
have higher implication in acoustic indices computation compared to
native species. Given that urbanisation gradient and non-native
dominated assemblage represents a combined effect. We expect that urban
sites have more saturated and complex soundscape patterns due to both
factors compared to deep forested sites.

\section{Methods}\label{methods}

\emph{Study area}

The study was conducted in Tahiti, French Polynesia. Tahiti is the
largest (1,042 km²) island of French Polynesia (Monnet et al. 1993). The
climate is tropical with an average annual temperature of 26,5 °C and
annual rainfall of 1700 mm in the coastal plain of the leeward coast of
Tahiti Nui, reaching 4000 mm on the windward coast and Tahiti Iti, and
up to 8000 mm in the center of Tahiti Nui (Laurent source quoin atlas
meteorologique accès?). Islands of the Pacific ocean undergoes strong
impoverishment with numerous bird extinctions due to first Polynesian
arrival in the 13th century and Europeans in the 18th century (Thibault
1988).(Bellwood 1978; Monnet et al. 1993; Steadman 1995) From 1960 to
1990, strong habitat modification took places due to development, costal
urbanisation and forest clearance (Blanchet 1986)(Monnet et al. 1993).
In Tahiti \enquote{\ldots{}} terrestrials bird species occurs, 12 of
them are non-native (Thibault et Cibois 2017). Introduced birds have
generally the highest densities in lowland, non-forested habitats such
as urban areas or agricultural fields, and lower densities in
high-elevation, contiguous forest, where native birds resides (Monnet et
al. 1993) (besoin source moins loin pour Tahiti or personal
observation).

\emph{Acoustic data collection}

Passive acoustic data was collected using SongMeter SM2 and SM4 digital
recording devices (Wildlife Acoustics, Maynard, MA, USA) from 2024 to
2026. In total 69 sites from different habitat types ( Forest = ,
\ldots{} rural or suburban = , .. Urban = ) were surveyed in valleys
along an attitudinal gradients. Stations recorded along the same
valley's gradient, so each gradient as staggering dates. Acoustic
recording units (ARUs) were programmed to record continuously and placed
at 1,5 meters above the ground.Devices recorded for \enquote{\ldots{}}
per station, with a total of \enquote{\ldots{}} hours recorded.
Recordings were made as .WAV files with a \enquote{\ldots{}} Hz sampling
rate and gain of \enquote{\ldots{}} dB, with monophonic and 16 Bits
formats.

\emph{Soundscape index metrics}

We selected nine relevant acoustic indices to encompass the sonic
emissions heterogeneity for the acoustic habitat and the soundtope of
each sampling site. Choice of indices were based on their relevance to
reflect patterns of variability, heterogeneity, complexity of both
soundscape level. The use of several acoustic indices was prioritize as
multi index analysis have been shown to be more effective than single
index investigations (Bradfer-Lawrence et al. 2020). Acoustic indices
were estimated for each 10-min recordings. We used indices median output
per hour for better control of index variation (Bradfer-Lawrence et al.
2019). Acoustic indices output were normalized to facilitate
comparisons.

Two different frequency limits were applied for the upper limit and a
low frequency threshold for devices glitches. For the soundtope the
limit was (\ldots{} Hz) to reflect Tahiti bird community emissions.
Recordings containing rainfall sounds were excluded from the calculation
at this soundscape level. For the acoustic habitat, the limit was
(\ldots{} Hz) to encompass all sonic emissions that could affect the
bird assemblages. Recording containing rainfall sounds were included in
the calculation at this soundscape level. All indices were estimated
using the default settings with the packages `soundecology' (\ldots) and
seewave (\ldots) packages in R (..).

(Table used for soundscape and bird acoustic community monitoring, with
hypothesis)

\emph{Habitat metrics}

Two metrics, the Hard Surface Area (HSA) and the open herbaceous cover
(HERB) were chosen to reflect habitat types and heterogeneity along the
attitudinal gradients. HSA and HERB ranges between 0 and 100 \%.
Respectively, higher values indicates the important cover of
human-modified surfaces and open herbaceous areas. Habitat covers map
was estimated using polygon surfaces in 100 meter radius circle around
each ARUs positions in Google Earth (Google Earth Pro.7.3.7.1327).

\emph{Avian caracterisation}

Two ten minute aural and visual identification in acoustic recordings
was conducted for each site, to obtain bird richness count. Aural and
visual acoustic identification were conducted from 06:00 am to 06:10 am
in recordings free of rainfall or strong acoustic disturbance. From the
aural identification database, several assemblage metrics were derived:
total richness (count), non-native richness (count), presence of native
bird (binomial), presence of Myna (binomial). Non-native bird species
encountered were \enquote{\ldots{}} Native bird species monitored were
\enquote{\ldots{}}.

\emph{Bird assemblage metrics in response to habitat characteristic
analysis}

Habitat as direct driver of bird assemblage was investigated. The three
bird assemblage metrics were modeled independently as responses
variables using \enquote{\ldots{}} for or \enquote{\ldots{}} for
\enquote{\ldots{}} distributions using the habitat metrics and altitude
as independent fixed variables. All generalized linear mixed-effects
models (GLMER) were made using the \enquote{GlmmTMB} package (Brooks et
al. 2017).

\emph{Acoustic indices in response to bird assemblage and habitat
metrics analysis}

Prior to analysis, acoustic indices output were tested for correlations.

To investigate if acoustic indices reflect metrics of native/non native
bird assemblage, GLMER were used to model each acoustic index (\ldots)
individually as response variable, for both acoustic habitat and
soundtope level. In response to the following independent variables: one
assemblage metric at a time and habitat metrics. All models included a
control fixed effect for the time (hour of the day) to account for
temporal autocorrelation and random fixed effects for the site and
commune. Models used a \enquote{\ldots{}} distribution with a
\enquote{\ldots{}} function. All statistical analyses were performed in
R v.().

\section*{References}\label{bibliography}
\addcontentsline{toc}{section}{References}

\protect\phantomsection\label{refs}
\begin{CSLReferences}{1}{1}
\bibitem[\citeproctext]{ref-bellwood1978}
Bellwood, Peter. 1978. \emph{The Polynesians: Prehistory of an Island
People}. Ancient Peoples et Places. Thames; Hudson.
\url{https://catalogue.nla.gov.au/catalog/2544013}.

\bibitem[\citeproctext]{ref-blackburn2001}
Blackburn, Tim M., et Richard P. Duncan. 2001. {«~Determinants of
Establishment Success in Introduced Birds~»}. \emph{Nature} 414 (6860):
195‑97. \url{https://doi.org/10.1038/35102557}.

\bibitem[\citeproctext]{ref-blackburn2009}
Blackburn, Tim M., Julie L. Lockwood, et Phillip Cassey. 2009.
\emph{Avian Invasions: The Ecology and Evolution of Exotic Birds}.
Oxford University Press.
\url{https://doi.org/10.1093/acprof:oso/9780199232543.001.0001}.

\bibitem[\citeproctext]{ref-bleach2015a}
Bleach, Iris, Christa Beckmann, Camila Both, Gregory Brown, et Richard
Shine. 2015. {«~Noisy neighbours at the frog pond: effects of invasive
cane toads on the calling behaviour of native Australian frogs~»}.
\emph{Behavioral Ecology and Sociobiology} 69 (février).
\url{https://doi.org/10.1007/s00265-015-1879-z}.

\bibitem[\citeproctext]{ref-boyer2014}
Boyer, Alison G., et Walter Jetz. 2014. {«~Extinctions and the Loss of
Ecological Function in Island Bird Communities~»}. \emph{Global Ecology
and Biogeography} 23 (6): 679‑88.
\url{https://doi.org/10.1111/geb.12147}.

\bibitem[\citeproctext]{ref-bradfer-lawrence2020}
Bradfer-Lawrence, Tom, Nils Bunnefeld, Nick Gardner, Stephen G. Willis,
et Daisy H. Dent. 2020. {«~Rapid assessment of avian species richness
and abundance using acoustic indices~»}. \emph{Ecological Indicators}
115 (août): 106400. \url{https://doi.org/10.1016/j.ecolind.2020.106400}.

\bibitem[\citeproctext]{ref-bradfer-lawrence2019}
Bradfer-Lawrence, Tom, Nick Gardner, Lynsey Bunnefeld, Nils Bunnefeld,
Stephen G. Willis, et Daisy H. Dent. 2019. {«~Guidelines for the Use of
Acoustic Indices in Environmental Research~»}. \emph{Methods in Ecology
and Evolution} 10 (10): 1796‑807.
\url{https://doi.org/10.1111/2041-210X.13254}.

\bibitem[\citeproctext]{ref-brooks2017}
Brooks, Mollie E., Kasper Kristensen, Koen J. van Benthem, et al. 2017.
{«~The R Journal: glmmTMB Balances Speed and Flexibility Among Packages
for Zero-inflated Generalized Linear Mixed Modeling~»}. \emph{The R
Journal} 9 (2): 378‑400. \url{https://doi.org/10.32614/RJ-2017-066}.

\bibitem[\citeproctext]{ref-davis2003}
Davis, Mark A. 2003. {«~Biotic Globalization: Does Competition from
Introduced Species Threaten Biodiversity?~»} \emph{BioScience} 53 (5):
481‑89.
\url{https://doi.org/10.1641/0006-3568(2003)053\%5B0481:BGDCFI\%5D2.0.CO;2}.

\bibitem[\citeproctext]{ref-druxf6ge2021}
Dröge, Saskia, Dominic Andreas Martin, Rouvah Andriafanomezantsoa, et
al. 2021. {«~Listening to a changing landscape: Acoustic indices reflect
bird species richness and plot-scale vegetation structure across
different land-use types in north-eastern Madagascar~»}.
\emph{Ecological Indicators} 120 (janvier): 106929.
\url{https://doi.org/10.1016/j.ecolind.2020.106929}.

\bibitem[\citeproctext]{ref-dyer2017}
Dyer, Ellie E., Phillip Cassey, David W. Redding, et al. 2017. {«~The
Global Distribution and Drivers of Alien Bird Species Richness~»}.
\emph{PLOS Biology} 15 (1): e2000942.
\url{https://doi.org/10.1371/journal.pbio.2000942}.

\bibitem[\citeproctext]{ref-eldridge2018}
Eldridge, Alice, Patrice Guyot, Paola Moscoso, Alison Johnston, Ying
Eyre-Walker, et Mika Peck. 2018. {«~Sounding out ecoacoustic metrics:
Avian species richness is predicted by acoustic indices in temperate but
not tropical habitats~»}. \emph{Ecological Indicators} 95 (décembre):
939‑52. \url{https://doi.org/10.1016/j.ecolind.2018.06.012}.

\bibitem[\citeproctext]{ref-fairbrass2017}
Fairbrass, Alison J., Peter Rennert, Carol Williams, Helena Titheridge,
et Kate E. Jones. 2017. {«~Biases of acoustic indices measuring
biodiversity in urban areas~»}. \emph{Ecological Indicators} 83
(décembre): 169‑77. \url{https://doi.org/10.1016/j.ecolind.2017.07.064}.

\bibitem[\citeproctext]{ref-farina2014}
Farina, Almo. 2014. \emph{Soundscape Ecology: Principles, Patterns,
Methods and Applications}. Springer Netherlands.
\url{https://doi.org/10.1007/978-94-007-7374-5}.

\bibitem[\citeproctext]{ref-farina2016}
Farina, Almo, et Philip James. 2016. {«~The acoustic communities:
Definition, description and ecological role~»}. \emph{Bio Systems} 147
(septembre): 11‑20.
\url{https://doi.org/10.1016/j.biosystems.2016.05.011}.

\bibitem[\citeproctext]{ref-farina2012}
Farina, Almo, et Nadia Pieretti. 2012. {«~The soundscape ecology: A new
frontier of landscape research and its application to islands and
coastal systems~»}. \emph{Journal of Marine and Island Cultures} 1
(juin): 21‑26. \url{https://doi.org/10.1016/j.imic.2012.04.002}.

\bibitem[\citeproctext]{ref-farina2013}
Farina, Almo, Nadia Pieretti, et Niki Morganti. 2013. {«~Acoustic
Patterns of an Invasive Species: The Red-Billed Leiothrix (
{\emph{Leiothrix Lutea}} Scopoli 1786) in a Mediterranean Shrubland~»}.
\emph{Bioacoustics} 22 (3): 175‑94.
\url{https://doi.org/10.1080/09524622.2012.761571}.

\bibitem[\citeproctext]{ref-fuller2015}
Fuller, Susan, Anne C. Axel, David Tucker, et Stuart H. Gage. 2015.
{«~Connecting soundscape to landscape: Which acoustic index best
describes landscape configuration?~»} \emph{Ecological Indicators} 58
(novembre): 207‑15. \url{https://doi.org/10.1016/j.ecolind.2015.05.057}.

\bibitem[\citeproctext]{ref-gage2014}
Gage, Stuart H., et Anne C. Axel. 2014. {«~Visualization of temporal
change in soundscape power of a Michigan lake habitat over a 4-year
period~»}. \emph{Ecological Informatics}, Ecological Acoustics, vol. 21
(mai): 100‑109. \url{https://doi.org/10.1016/j.ecoinf.2013.11.004}.

\bibitem[\citeproctext]{ref-gasc2018}
Gasc, A., J. Anso, J. Sueur, H. Jourdan, et L. Desutter-Grandcolas.
2018. {«~Cricket Calling Communities as an Indicator of the Invasive Ant
Wasmannia Auropunctata in an Insular Biodiversity Hotspot~»}.
\emph{Biological Invasions} 20 (5): 1099‑111.
\url{https://doi.org/10.1007/s10530-017-1612-0}.

\bibitem[\citeproctext]{ref-halfwerk2011}
Halfwerk, Wouter, Leonard J. M. Holleman, C(Kate). M. Lessells, et Hans
Slabbekoorn. 2011. {«~Negative Impact of Traffic Noise on Avian
Reproductive Success~»}. \emph{Journal of Applied Ecology} 48 (1):
210‑19. \url{https://doi.org/10.1111/j.1365-2664.2010.01914.x}.

\bibitem[\citeproctext]{ref-hopkins2022}
Hopkins, Jaimie M., Will Edwards, et Lin Schwarzkopf. 2022. {«~Invading
the Soundscape: Exploring the Effects of Invasive Species{'} Calls on
Acoustic Signals of Native Wildlife~»}. \emph{Biological Invasions} 24
(11): 3381‑93. \url{https://doi.org/10.1007/s10530-022-02856-w}.

\bibitem[\citeproctext]{ref-hunt2024}
Hunt, Noah J., Thomas Ibanez, Adam A. Pack, et Patrick J. Hart. 2024.
{«~Signal Partitioning Between Native and Introduced Forest Birds of
Hawai{`}i Island~»}. \emph{Frontiers in Ecology and Evolution} 12
(août). \url{https://doi.org/10.3389/fevo.2024.1399455}.

\bibitem[\citeproctext]{ref-lellouch2014}
Lellouch, Laurent, Sandrine Pavoine, Frédéric Jiguet, Hervé Glotin, et
Jérôme Sueur. 2014. {«~Monitoring Temporal Change of Bird Communities
with Dissimilarity Acoustic Indices~»}. \emph{Methods in Ecology and
Evolution} 5 (6): 495‑505.
\url{https://doi.org/10.1111/2041-210X.12178}.

\bibitem[\citeproctext]{ref-lohr2003}
Lohr, Bernard, Timothy F. Wright, et Robert J. Dooling. 2003.
{«~Detection and discrimination of natural calls in masking noise by
birds: estimating the active space of a signal~»}. \emph{Animal
Behaviour} 65 (4): 763‑77. \url{https://doi.org/10.1006/anbe.2003.2093}.

\bibitem[\citeproctext]{ref-mcgill2006}
McGill, Brian J., Brian J. Enquist, Evan Weiher, et Mark Westoby. 2006.
{«~Rebuilding community ecology from functional traits~»}. \emph{Trends
in Ecology \& Evolution} 21 (4): 178‑85.
\url{https://doi.org/10.1016/j.tree.2006.02.002}.

\bibitem[\citeproctext]{ref-medeiros2017}
Medeiros, Camila Ineu, Camila Both, Taran Grant, et Sandra Maria Hartz.
2017. {«~Invasion of the Acoustic Niche: Variable Responses by Native
Species to Invasive American Bullfrog Calls~»}. \emph{Biological
Invasions} 19 (2): 675‑90.
\url{https://doi.org/10.1007/s10530-016-1327-7}.

\bibitem[\citeproctext]{ref-merchant2015}
Merchant, Nathan D., Kurt M. Fristrup, Mark P. Johnson, et al. 2015.
{«~Measuring Acoustic Habitats~»}. \emph{Methods in Ecology and
Evolution} 6 (3): 257‑65. \url{https://doi.org/10.1111/2041-210X.12330}.

\bibitem[\citeproctext]{ref-monnet1993}
Monnet, Claude, Jean-Claude Thibault, et Albert Varney. 1993.
{«~Stability and Changes During the Twentieth Century in the Breeding
Landbirds of Tahiti (Polynesia)~»}. \emph{Bird Conservation
International} 3 (4): 261‑80.
\url{https://doi.org/10.1017/S0959270900002550}.

\bibitem[\citeproctext]{ref-pijanowski2011}
Pijanowski, Bryan C., Luis J. Villanueva-Rivera, Sarah L. Dumyahn, et
al. 2011. {«~Soundscape Ecology: The Science of Sound in the
Landscape~»}. \emph{BioScience} 61 (3): 203‑16.
\url{https://doi.org/10.1525/bio.2011.61.3.6}.

\bibitem[\citeproctext]{ref-ribeiro2022}
Ribeiro, José W., Kristopher Harmon, Gabriel Augusto Leite, Tomaz
Nascimento de Melo, Jack LeBien, et Marconi Campos-Cerqueira. 2022.
{«~Passive Acoustic Monitoring as a Tool to Investigate the Spatial
Distribution of Invasive Alien Species~»}. \emph{Remote Sensing} 14
(18): 4565. \url{https://doi.org/10.3390/rs14184565}.

\bibitem[\citeproctext]{ref-robert2019}
Robert, Aloïs, Thierry Lengagne, Martim Melo, et al. 2019. {«~The Theory
of Island Biogeography and Soundscapes: Species Diversity and the
Organization of Acoustic Communities~»}. \emph{Journal of Biogeography}
46 (9): 1901‑11. \url{https://doi.org/10.1111/jbi.13611}.

\bibitem[\citeproctext]{ref-rogers2021}
Rogers, Andrew M, Andrea S Griffin, Françoise Lermite, Berndt van
Rensburg, Carla Archibald, et Salit Kark. 2021. {«~The role of invasion
and urbanization gradients in shaping avian community composition~»}.
\emph{Journal of Urban Ecology} 7 (1): juab030.
\url{https://doi.org/10.1093/jue/juab030}.

\bibitem[\citeproctext]{ref-rudge2026}
Rudge, Leah, Angela Holland, John Schmit, et Greg Shriver. 2026. {«~In
urban national parks acoustic indices represent avian habitat quality
better than avian diversity~»}. \emph{Environmental and Sustainability
Indicators} 30 (février): 101174.
\url{https://doi.org/10.1016/j.indic.2026.101174}.

\bibitem[\citeproctext]{ref-russell2019}
Russell, James C., et Christoph Kueffer. 2019. {«~Island Biodiversity in
the Anthropocene~»}. \emph{Annual Review of Environment and Resources}
44: 31‑60. \url{https://doi.org/10.1146/annurev-environ-101718-033245}.

\bibitem[\citeproctext]{ref-suxe1nchez-giraldo2021}
Sánchez-Giraldo, Camilo, Camilo Correa Ayram, et Juan M. Daza. 2021.
{«~Environmental sound as a mirror of landscape ecological integrity in
monitoring programs~»}. \emph{Perspectives in Ecology and Conservation}
19 (3): 319‑28. \url{https://doi.org/10.1016/j.pecon.2021.04.003}.

\bibitem[\citeproctext]{ref-sax2007}
Sax, Dov F., John J. Stachowicz, James H. Brown, et al. 2007.
{«~Ecological and evolutionary insights from species invasions~»}.
\emph{Trends in Ecology \& Evolution} 22 (9): 465‑71.
\url{https://doi.org/10.1016/j.tree.2007.06.009}.

\bibitem[\citeproctext]{ref-sayol2021}
Sayol, Ferran, Robert S. C. Cooke, Alex L. Pigot, et al. 2021. {«~Loss
of functional diversity through anthropogenic extinctions of island
birds is not offset by biotic invasions~»}. \emph{Science Advances} 7
(46): eabj5790. \url{https://doi.org/10.1126/sciadv.abj5790}.

\bibitem[\citeproctext]{ref-seebens2017}
Seebens, Hanno, Tim M. Blackburn, Ellie E. Dyer, et al. 2017. {«~No
Saturation in the Accumulation of Alien Species Worldwide~»}.
\emph{Nature Communications} 8 (1): 14435.
\url{https://doi.org/10.1038/ncomms14435}.

\bibitem[\citeproctext]{ref-shannon2016}
Shannon, Graeme, Megan F. McKenna, Lisa M. Angeloni, et al. 2016. {«~A
Synthesis of Two Decades of Research Documenting the Effects of Noise on
Wildlife~»}. \emph{Biological Reviews} 91 (4): 982‑1005.
\url{https://doi.org/10.1111/brv.12207}.

\bibitem[\citeproctext]{ref-soares2021}
Soares, Filipa C., Ana I. Leal, Jorge M. Palmeirim, et Ricardo F. De
Lima. 2021. {«~Niche Differences May Reduce Susceptibility to
Competition Between Native and Non{-}Native Birds in Oceanic Islands~»}.
\emph{Diversity and Distributions} 27 (8): 1507‑18.
\url{https://doi.org/10.1111/ddi.13298}.

\bibitem[\citeproctext]{ref-soares2022}
Soares, Filipa, Ricardo Lima, Jorge Palmeirim, Pedro Cardoso, et Ana
Rodrigues. 2022. {«~Combined effects of bird extinctions and
introductions in oceanic islands: Decreased functional diversity despite
increased species richness~»}. \emph{Global Ecology and Biogeography} 31
(mars): 1172‑83. \url{https://doi.org/10.1111/geb.13494}.

\bibitem[\citeproctext]{ref-sobral2016}
Sobral, Fernando L., Alexander C. Lees, et Marcus V. Cianciaruso. 2016.
{«~Introductions Do Not Compensate for Functional and Phylogenetic
Losses Following Extinctions in Insular Bird Assemblages~»}.
\emph{Ecology Letters} 19 (9): 1091‑100.
\url{https://doi.org/10.1111/ele.12646}.

\bibitem[\citeproctext]{ref-spatz2017}
Spatz, Dena R., Kelly M. Zilliacus, Nick D. Holmes, et al. 2017.
{«~Globally threatened vertebrates on islands with invasive species~»}.
\emph{Science Advances} 3 (10): e1603080.
\url{https://doi.org/10.1126/sciadv.1603080}.

\bibitem[\citeproctext]{ref-steadman1995}
Steadman, David W. 1995. {«~Prehistoric Extinctions of Pacific Island
Birds: Biodiversity Meets Zooarchaeology~»}. \emph{Science} 267 (5201):
1123‑31. \url{https://doi.org/10.1126/science.267.5201.1123}.

\bibitem[\citeproctext]{ref-thibault2017}
Thibault, Jean-Claude, et Alice Cibois. 2017. \emph{Birds of Eastern
Polynesia. A biogeographic Atlas}.

\bibitem[\citeproctext]{ref-towsey2014}
Towsey, Michael, Jason Wimmer, Ian Williamson, et Paul Roe. 2014. {«~The
use of acoustic indices to determine avian species richness in
audio-recordings of the environment~»}. \emph{Ecological Informatics},
Ecological Acoustics, vol. 21 (mai): 110‑19.
\url{https://doi.org/10.1016/j.ecoinf.2013.11.007}.

\bibitem[\citeproctext]{ref-villanueva-rivera2011}
Villanueva-Rivera, Luis J., Bryan C. Pijanowski, Jarrod Doucette, et
Burak Pekin. 2011. {«~A Primer of Acoustic Analysis for Landscape
Ecologists~»}. \emph{Landscape Ecology} 26 (9): 1233‑46.
\url{https://doi.org/10.1007/s10980-011-9636-9}.

\bibitem[\citeproctext]{ref-whittaker2017}
Whittaker, Robert J., José María Fernández-Palacios, Thomas J. Matthews,
Michael K. Borregaard, et Kostas A. Triantis. 2017. {«~Island
biogeography: Taking the long view of nature{'}s laboratories~»}.
\emph{Science} 357 (6354): eaam8326.
\url{https://doi.org/10.1126/science.aam8326}.

\end{CSLReferences}




\end{document}
